\section{Error analysis in least squares}


Let $(\X \times \Y, \mu)$ be a probability space, with data $\X \subseteq \R^d$, $\Y\subseteq\R$, and let $\ell:\Y\times\Y \to [0,\infty)$ be a measurable loss function.

For a continuous function of the data $f:\X\to\Y$, we give the following definition. 

\begin{definition}[Expected risk]
The expected risk of $f:\X\to\Y$ is 
\begin{equation}
    L(f) \coloneq \expv_{(x,y)\sim\mu} \big[\ell(y,f(x))\big] = \int_{\X\times\Y} \ell(y,f(x)) \de \mu(x,y)
\end{equation}
\end{definition}

In general, we aim at finding \begin{equation}\label{pr: exprisk}f_\mu \coloneq \argmin_{f\in\F} L(f)\end{equation}
where the minimum runs over the space $\F$ of functions for which $L(f)$ is well defined (the space of measurable functions). Given a learning algorithm and some data ${\{x_i,y_i\}}_{i=1}^n, x_i \in \X, y_i \in \Y \,\forall i$, we obtain a learning solution $\hat f$. Then, the natural error measure is the excess risk
\begin{definition}[Excess risk]
    The excess risk of a learning problem with target $f_\mu$ and learning solution $\hat f$ is
    \[ L(\hat f) - L(f_\mu)\]
    where $L(\cdot)$ is the expected risk.
\end{definition}

Since this minimum in (\ref{pr: exprisk}) can be hard or impossible to compute, we give the following definition. 

\begin{definition}[Empirical risk]
    Given a probability space $(\X \times \Y, \mu)$, a dataset ${\{x_i,y_i\}}_{i=1}^n, x_i \in \X, y_i \in \Y \,\forall i$, a continuous measurable function $f:\X\to\Y$, and a measurable loss function $\ell:\Y\times\Y \to [0,\infty)$, the empirical risk of $f$ is
    \begin{equation}\label{eq: emprisk}
        \hat L(f) \coloneq \sum_{i=1}^n \ell(f(x_i), y_i)
    \end{equation}
    and it is a finite-sample, tractable version of the expected risk. 
\end{definition}
The problem becomes to find \begin{equation} \min_{f} \hat L(f)\end{equation}
which is known as {\it empirical risk minimization (ERM)},
where $\mathcal{H}$ is the hypotheses set, a restricted set of possible solutions that we choose based on our problem. 
% Our learning algorithm will now give $\hat f = \argmin_{\f \in \cH} \hat L (f)$.
\\

Now, we focus our analysis on regularized least squares. We consider a linear combination of the features
\[
f(x) = \langle w,x \rangle, \quad w\in\R^d
\]
hence we have $\cH =\{f:\X\to\Y \,\vert\, f(x) = \langle w,x \rangle \text{ for some } w \in\R^d\}$;
the squared loss
\[
\ell\big(f(x),y\big) = {\big(y - \langle w,x\rangle\big)}^2
\]
with associated expected risk
\[
L(w) = \int_{\X\times\Y} \big(y - \langle w,x\rangle\big)^2\de\mu(x,y)
\]
(indicated as $L(w)$ with an abuse of notation).
With a finite dataset of size $n$, our estimator will be 
\begin{equation}\label{eq: estim}
    \hat w_\lambda = \argmin_{w\in\R^d} \frac{1}{n} \sum_{i=1}^n \big(y_i - \langle x_i,w\rangle\big)^2 + \lambda \Vert w \Vert^2
\end{equation}
minimizing the empirical risk plus a regularization term with $\lambda >0$, and the natural error measure will be
\begin{equation}\label{eq: errmeasureLS}
    L(\hat w_\lambda) - \min_{w\in\R^d} L(w)
\end{equation}
i.e. the excess risk, which is a positive random variable. We are interested in bounding
\begin{equation}\label{pr: eps1}
    \Pr \bigg[L(\hat w_\lambda) - \min_{w\in\R^d} L(w) \le \varepsilon\bigg], \quad \forall\varepsilon>0
\end{equation}

its cumulative distribution. 

Alternatively, we assume $\exists \, \kappa,M>0$ such that $\Vert x \Vert \le \kappa$ and $\vert y \vert \le M$ almost surely, and that there is a $w^* \in \R^d$ that minimizes $L(w)$. Then, we look at the error measure
\[
\Vert \hat w_\lambda - w^* \Vert
\]
and consider 
\begin{equation}\label{pr: eps2}
    \Pr \bigg[\Vert \hat w_\lambda - w^* \Vert \le \varepsilon\bigg]
\end{equation}
which is easier to analyse than (\ref{pr: eps1}), but provides worse bounds (any bound on (\ref{pr: eps2}) implies bounds on (\ref{pr: eps1})). 

\paragraph{\bf Optimality conditions and other properties}

We consider our problem under 3 different lenses, going from the infinite data case with no regularization to the empirical case that we actually implement, and provide optimality conditions and other results for each of them. 
\\

We start by considering problem
\[
\min_{w\in\R^d} L(w)
\]
and we define the following useful quantities
\begin{gather*}
    \Sigma:\R^d\to\R^d, \qquad \Sigma w = \int_{\X\times\Y} xx^tw \de \mu(x,y) = \int_{\X} xx^tw \de \mu_\X(x)\\
    h\in\R^d, \qquad h= \int_{\X\times\Y} xy\de\mu(x,y)
    \end{gather*}
        where $\mu_\X (x) = \int_{\Y} \mu(x,y) \de y$ is the marginal distribution of $x$. 
\begin{proposition}[Infinite data, no regularization]\label{prop: infinite,noreg}
Consider the minimization problem 
\[
\min_{w\in\R^d} L(w)
%= \int_{\X\times\Y} (y-\langle w,x \rangle ) ^ 2 \de \mu(x,y)
\]
%for $\mu$ probability over the data space $\X\times\Y$ and $\ell(x,y) = (y-\langle w,x \rangle ) ^ 2$ measurable loss function.
then, the following hold:
    \begin{enumerate}[label=(\roman*)]
        \item \[\nabla L(w)=0 \quad\Longleftrightarrow\quad \Sigma w = h\]
        with $\Sigma, h$ as above;
        \item in particular, \[w^* = \Sigma^{-1}h\]
        where $w^*$ minimizes $L(w)$;
        \item additionally, we have
        \[L(w) - L(w^*) = \langle w-w^*, \Sigma(w-w^*)\rangle\]
    \end{enumerate}
    \begin{proof}
        The proof of item $(i)$ follows by expanding the gradient of $L(w)$ and noticing that the two conditions coincide:
        \begin{align*}
            \nabla L(w)& =\frac{\partial}{\partial w} \int (y-\langle w,x \rangle)^2 \de\mu = \int \frac{\partial}{\partial w} (y-\langle w,x \rangle)^2 \de\mu \\
            & = \int 2 (y-\langle w,x \rangle) (-x) \de \mu = 2 \int - yx + \langle w,x \rangle x \de\mu;\\
            \nabla L(w)=0&\Leftrightarrow \int yx \de \mu = \int x\langle w,x \rangle \de \mu \Leftrightarrow h = \Sigma w
        \end{align*}
        where the first line follows from continuity of $\ell(x,y)=(y-\langle w,x \rangle)^2$. Item $(ii)$ follows directly by definition of minimum and item $(i)$. 
        Lastly, item $(iii)$ follows by manipulating the expression of $L(w)$:

        \begin{align*}
            L(w) &= \int (y-\langle w,x\rangle) ^ 2\de \mu = \int (y - \langle w^*,x\rangle + \langle w^*,x \rangle - \langle w,x\rangle)^2 \de \mu\\
            &= \int (y-\langle w^*,x\rangle)^2 \de \mu + \int (\langle w^*,x\rangle-\langle w,x\rangle)^2 \de \mu - 2\int (y - \langle w^*,x\rangle)(\langle w^*,x\rangle-\langle w,x\rangle) \de \mu \\
            & = L(w^*) + \int (\langle w^*,x\rangle-\langle w,x\rangle)^2 \de \mu - 2\int (y - \langle w^*,x\rangle)(\langle w^*,x\rangle-\langle w,x\rangle) \de \mu
        \end{align*}
        the latter term is
        \begin{align*}
            \int (y-\langle w^*,x\rangle) (\langle x,w^*-w\rangle) \de \mu = \bigg\langle \int (y-\langle w^*,x\rangle)x\de mu, w^*-w\bigg\rangle = \langle h-\Sigma w^*, w^*-w\rangle = 0
        \end{align*}
        where we used the property (\ref{app: trick}) of the integral and continuous functions, and item $(ii)$. Hence,
        \begin{align*}
            L(w) - L(w^*) &= \int (\langle w^*,x\rangle-\langle w,x\rangle)^2 \de \mu  = \int \langle w^*-w,x\rangle\langle w^*-w,x\rangle \de \mu = \bigg\langle w^* - w, \int \langle w^* -w, x\rangle x \de \mu \bigg\rangle \\
            & = \langle w^*-w , \Sigma (w^*-w)\rangle
        \end{align*}
        
    \end{proof}
\end{proposition}

Similarly, we have the following result.
\begin{proposition}[Infinite data, with regularization]\label{prop: infinite,reg}
    Consider the problem 
    \[
    \min_{w\in\R^d} L_\lambda(w) = L(w) + \lambda \Vert w\Vert^2
    \]
    for some $\lambda >0$. Then, the following hold:
    \begin{enumerate}[label=(\roman*)]
        \item the corresponding optimality condition is
        \[\nabla L_\lambda(w) = 0 \quad \Longleftrightarrow \quad (\Sigma + I\lambda)w = h\]
        where $I:\R^d \to \R^d$ is the identity map;
        \item in particular, 
        \[ w_\lambda = (\Sigma + \lambda I)^{-1} h\]
        where $w_\lambda$ minimizes $L_\lambda(w)$.
    \end{enumerate}
    \begin{proof}
        The proof follows by the same arguments as those in proposition (\ref{prop: infinite,noreg}). 
    \end{proof}
\end{proposition}

Now, we turn to analysing the case of empirical least squares, and define the associated quantities
\begin{gather*}
    \hat\Sigma:\R^d\to\R^d, \qquad \hat \Sigma = \frac{1}{n}\sum_{i=1}^n x_i x_i^t\\
    \hat h\in\R^d, \qquad \hat h= \frac{1}{n}\sum_{i=1}^n x_iy_i
    \end{gather*}

We consider the problem
\[\min_{w\in\R^d} \hat L_\lambda\]
with empirical risk
\[\hat L_\lambda = \sum_{i=1}^n (y_i - \langle w, x_i\rangle)^2 + \lambda \Vert w \Vert ^2, \quad \lambda>0\]
(which follows from definition (\ref{eq: emprisk}) and adding a regularizer as in $L_\lambda$).
Repeating the same arguments, we obtain
\[\hat w_\lambda = (\hat \Sigma + \lambda I)^{-1} \hat h\]
where $\hat w_\lambda$ minimizes $\hat L_\lambda$.
